
%%%%%%%%%%%%%%%%%%%%%%%%%%%%%%%%%%%%%%%%%%%%%%%%%%%%%%%%%%%%

\chapter*{Foreword}
\phantomsection \addcontentsline{toc}{chapter}{Foreword}
 % telling stories is important for knowledge transmission

% consequently, le ton sera variable, non always strictly scientific

% Cette thèse rend compte de ce qui m'a occupé pendant les trois dernières années. Il ne s'agit donc pas exclusivement de mon travail de recherche mais aussi, dans une moindre mesure, de ce que j'ai enseigné aux étudiants, transmis au publique, ainsi que ce que j'ai appris, et tout ce qui a nourris mes pensées et conversations avec mon entourage.
% Le chapitre premier, et notamment son commencement, prend donc la forme d'une large conduite qui rétrécit doucement vers le cœur de mon sujet de recherche. 
% La première hypothèse de cette thèse est que, garder constamment en tête un point de vue large sur l'état du climat est nécessaire dans l'étude de n'importe laquelle de ces composantes, aussi spécifique que puisse être le sujet.
% Le lecteur ne devra pas alors s'étonner de ne parler nuage et microphysique qu'après une quinzaine de pages. 

% Le manuscrit est rédigé en anglais qui n'est, de toute évidence, pas ma langue maternelle. Certaines de mes idées pourront alors paraître confuses, ou d'expression maladroite, bien que je me sois efforcer de lisser autant que faire se peut ces variations de rythme. Malgré tout cela, j'espère que le lecteur sentira, dans son parcours à travers l'ouvrage, quelque peu du plaisir que j'ai eu à l'écrire.

%%The first intention of this thesis is that having in mind a broad view on the climate state is necessary for the study of any of its components, as specialised as it could be.

This thesis accounts for what has occupied me over the past three years, bla bla blahh 
